\setlength{\parindent}{0em}

\chapter{Desarrollo}\label{cap3}\thispagestyle{fancy}

\vspace{-1.3cm}

\section{Desarrollo del autómata}

El desarrollo del sistema se divide en diferentes fases dentro de la ETL, destacando la carga de datos históricos, su limpieza y filtrado,
su transformación en las diferentes características para alimentar a los modelos de aprendizaje automático.

% Además, están las partes de 

% Se quiere usar una arquitectura de medallón, para representar la utilidad de los datos, en primer lugar están los datos en crudo que se cargan

El objetivo es poder generalizar el código, que no importe de qué activos sean los datos el flujo de datos puede reutilizarse.

Para lograr este objetivo he decidido implementar la lógica del código dentro de un paquete instalable de python usando
\mintinline{bash}{pip install -e} y he decidido usar la siguiente estructura organizativa.

\dirtree{%
    .1 /.
    .2 paths.py.
    .2 get\_info.
    .2 data.
        .3 loader.py.
        .3 cleaner.py.
        .3 resampler.py.
        .3 to\_parquet.py.
    .2 features.
    .2 supervised.
    .2 unsupervised.
    .2 portfolio.
    .2 risk.
    .2 execution.
    .2 logging.
    .2 retrain.
}

Hay funciones que están en algunos de los directorios que pueden servir para varias entapas del proceso, principalmente aquellas encargadas de
obtener los datos desde los archivos.

A continuación paso a explicar cada una de las fases del flujo de los datos y las funciones utilizadas en ellas.

\subsection{Carga y limpieza}

Se me han facilitado los datos históricos OHLCV de distintos futuros con los que voy a operar.

En un primer análisis me doy cuenta de que los datos temporales no están regularizados al formato UTC, por lo que es el primer cambio que realizo.
Usando de referencia las horas de apertura y cierre de cada futuro, agrupo el volumen por horas para cada activo puedo detectar las horas con poca
o nula actividad, con ello puedo inferir la zona horaría usada.

\todo{Explicar mejor cómo funciona la función de inferencia}

Después, decido filtrar los datos que claramente no deberían de estar, estos son los que aparecen en los fines de semana. Los pedidos realoizados en
las horas de descanso las mantengo por que pueden haber sido retrasos en la ejecución por lo que son operaciones reales no realizadas en el momento
indicado.

Una vez estandarizado el tiempo aprovecho la función de inferencia y añado una priemra columna derivada que indica el día de \textit{trading} aunque
al ser una fecha no se usará directamente para el entrenamiento si me ayudará para calcular valores acumulados diarios.

Por otra parte como mi intención es obtener métricas de comparación entre activos decido recortar el número de filas disponibles en los diferentes
archivos para tener un conjunto homogéneo. Además para mejorar el rendimiento y disminuir el tamaño de los archivos convierto los datos que estaban
en formato csv a parquet.

\subsection{Cálculo de características}

Una vez con los datos OHLCV limpios, llega el momento de aplicar diferentes funciones de diversas librerías para obtener características derivadas
que servirán para alimentar a los modelos.

Hay características de diferentes tipos, momento, estructural, temporal, de tendencia, de volatilidad, de volumen y características derivadas de
otras. Al ser una estratégia multidía también he hecho uso de los swing charts y de características derivadas de ellas, además de otras
características obtenidas de comparar diferentes activos

Para obtener estas características primero agrupo los datos que venían minuto a minuto a agrupaciones de 4 horas, de esto modo se evita tener el
ruido de señales tan cortas y es un buen número para una estratégia interdía donde una posición puede estar abierta varios días

\subsection{Selección de características}

Una vez calculadas todas las características es necesario reducir ese número porque es posible que muchas de ellas estén interrelacionadas

Para ello se aplican varias técnicas, empezando por la más básica, un filtro de alta correlación, además de eliminar las características que puedan
predecir el futuro, las no cíclicas o no estacionarias pueden confundir a los modelos indicando que en valor que quizás representaba el día de la
semana sea considerado como más importante a mayor número

También filtro usando el algorítmo de LASSO usando un modelo de aprendizaje automático clásico de sklearn que permité mayor explicabilidad

\subsection{Aprendizaje supervisado}

El uso de modelos de aprendizaje supervisado se centra en modelos de clasificación binaria, en este caso en dos etapas, y dentro de cada etapa se
utilizan dos modelos de aprendizaje profundo usando PyTorch, LSTM y GRU. Cada uno de los modelos utilizados tiene una función, LSTM ..., por su parte GRU.
Para ambas fases se realiza una votación equitativa, si la media supera el umbral indicado para cada fase se pasa a la siguiente
\notes{Probar a ver otros métodos de votación además de usar la media}

Existe una tercera etapa, en este caso de regresión que sigue la misma lógica, usando los dos modelos

\subsection{Aprendizaje no supervisado}

Pasando ahora a los modelos, primero voy a explicar las dos situaciones donde voy a implementar modelos de aprendizaje no supervisado

En primer lugar para agrupar los diferentes activos en diferentes grupos. En vez de realizar agrupaciones arbitrarias por ejemplo todas las monedas
juntas dejo que el sistema decida cuales son los activos más parecidos

\todo{Decidir el algoritmo a usar para UL}
% \review{Prueba cuadro \textbackslash review}

\section{Despliegue usando contenedores de Docker}

Una vez que se ha desarrollado el autómata y se ha comprobado que su funcionamiento es correcto, he decidido añadir el proyecto a un contenedor de
docker para facilitar el despligue en otros equipos

\subsection{Configuración del contenedor principal}

Por cómo he desarrollado el sistema y por velocidad todo va a estar incluido en un único contenedor.

No requiere demasidos cambios ya que voy a copiar los códigos desarrollados dentro del contenedor

\subsection{Bases de datos}

Al tener el proyecto dentro de un contenedor veo lógico tener una base de datos con diferentes tablas para guardar diferentes métricas

\subsection{Alertas y avisos}

Por último una vez que todo funciona en contenedores he decidido añadir un sistema de visualización de toda la información que se recopila en
directo y en diferido

Uso Prometheus, Grafana y Alertmanager este último conectado a otro servicio que se conecta a la API de whatsapp

\subsection{Orquestación y MLOps}

\notes{Tanto Airflow como MLflow, ver si hay tiempo para implementarlos}

Una vez desplegado un sistema robusto que cuenta con un autómata funcional y con alertas tanto de seguridad como para mostrar información, decido
implementar el uso de Airflow como orquestrador.

En general el uso de Airflow está muy ligado al uso de MLflow ya que será el uso prinicpal de esta herramienta por mi parte.

En MLflow 

% \subsection{MLOps}

% Uso de MLflow